\section{Тестирование и верификация}

\subsection{Модульные тесты}

Для верификации изолированных компонентов эмулятора разработан
набор модульных тестов на базе фреймворка GoogleTest. Тесты
организованы в четыре файла: \texttt{test\_executors.cpp}
(исполнители инструкций), \texttt{test\_barrier.cpp} (барьерная синхронизация),
\texttt{test\_coalescing.cpp} (объединение транзакций памяти)
и \texttt{test\_ptx\_asm.cpp} (ассемблирование PTX-фрагментов).

Разбивка тестов по файлам соответствует границам архитектурных компонентов:
каждый файл тестирует ровно один субмодуль, что упрощает локализацию
регрессий при внесении изменений в конкретный компонент.

Покрытие кода измеряется инструментом \texttt{lcov}.

Покрытие строк составляет \textbf{80.4\%}, покрытие ветвей~--- \textbf{78.7\%}.
Непокрытые ветви сосредоточены преимущественно в обработчиках ошибочных
состояний (\texttt{global\_context}, путь выделения GPU-памяти),
которые требуют наличия CUDA Runtime и не тестируются модульными тестами.
Модульные тесты покрывают изолированные компоненты и служат
страховочной сеткой при доработке кода, однако не заменяют интеграционное
тестирование на реальных ядрах.

Достигнутый уровень покрытия~--- свыше 80\% строк и ветвей~--- является
достаточным для верификации основного функционального пути,
включающего все классы инструкций и ключевые граничные случаи,
описанные в Главе~5. Для производственного инструмента это значение
следует повышать, в первую очередь за счёт добавления тестов,
имитирующих ошибочные сценарии инициализации контекста.

\subsection{Часть I --- Интеграционное тестирование CUDA-ядер}

\textbf{Методология.} Каждый интеграционный тест компилируется
с \texttt{nvcc} и запускается под эмулятором: \texttt{cuemu ./binary}.
Тест самостоятельно сравнивает результат с CPU-эталоном, возвращая
код возврата~0 при успехе. Для оценки производительности каждое
ядро запускается 10~раз; фиксируются медиана и стандартное отклонение
времени \texttt{KernelLaunch}. Тесты запускаются на машине без GPU
(процессор AMD Ryzen 7 3700U, 8~ГБ RAM).

Семь ядер выбраны таким образом, чтобы перекрывать качественно различные
вычислительные паттерны: от простейших поэлементных операций (\texttt{vadd})
через ядра с нетривиальным управлением потоком и разделяемой памятью
(\texttt{softmax}, \texttt{mmul}) до сложных алгоритмов со вложенными
барьерами и динамическим числом итераций (\texttt{attention}, \texttt{conj\_grad}).
Успешное прохождение всего набора свидетельствует о сквозной корректности
интеграции всех компонентов эмулятора.

\subsubsection{\texttt{vadd\_custom} --- поэлементное сложение векторов}

\textbf{Описание теста.}
Ядро складывает два массива \texttt{float} длиной $N=1024$ на одном
блоке из 32~потоков. CPU-эталон вычисляется скалярным циклом;
допуск~--- абсолютная погрешность $|gpu[i] - ref[i]| \leq 10^{-6}$.
Тест проверяет базовый путь: загрузка из глобальной памяти
(\texttt{ld.global.f32}), арифметику (\texttt{add.f32}),
запись в глобальную память (\texttt{st.global.f32}).

\textbf{Тестирование.}
Фиксируются: код возврата (0 = PASS), максимальная абсолютная
ошибка, число активных потоков на инструкцию
(ожидается ровно 32~--- дивергенции нет).

\textbf{Замеры производительности.}
Ядро содержит минимальное число инструкций (3 на поток),
что делает его базовой точкой отсчёта. Замер проводился
методом «50 повторений одного запуска ядра в одном процессе»~---
для исключения шума от запуска динамического загрузчика;
фиксировалось минимальное значение из трёх серий.
Машина: AMD Ryzen 7 3700U, 8~ГБ RAM, без GPU.

\begin{table}[h]
  \caption{Время исполнения \texttt{vadd\_custom} при различных $N$
           (медиана по 50 повторениям)}
  \label{tab:vadd_perf}
  \centering
  \begin{tabular}{rrrr}
    \toprule
    $N$ & Без проф., мс & Со сбором проф., мс & Накл. расходы, \% \\
    \midrule
    256   & 0.44 & 1.02 & 132 \\
    1024  & 1.74 & 3.89 & 124 \\
    4096  & 8.34 & 16.92 & 103 \\
    16384 & 39.69 & 74.50 & 88 \\
    \bottomrule
  \end{tabular}
\end{table}

\textbf{Вывод.}
Накладные расходы профилировщика составляют 88--132\%~--- то есть примерно
удваивают время выполнения. Относительные накладные расходы убывают с ростом $N$:
при малом числе инструкций записи в \texttt{ProfilingRecord}
сопоставимы с суммарным временем исполнения варпов;
при больших $N$ исполнение варпов начинает доминировать.
Данное поведение соответствует ожидаемой модели: накладные расходы профилировщика
линейны по числу инструкций, а не по $N$ напрямую.

Двукратное увеличение времени при включённом профилировании является
приемлемым для задач анализа и отладки: пользователь, обращающийся
к профилировщику, осознанно принимает это замедление в обмен на детальные
метрики, недоступные при обычном запуске.

\subsubsection{\texttt{gelu\_custom} --- функция активации GELU}

\textbf{Описание теста.}
Ядро применяет GELU-аппроксимацию:
\[
  \mathrm{GELU}(x) = x \cdot \tfrac{1}{2}
  \Bigl(1 + \tanh\!\Bigl(\sqrt{\tfrac{2}{\pi}} \cdot
  (x + 0.044715 x^3)\Bigr)\Bigr)
\]
к массиву из $N=4096$ значений. CPU-эталон~--- \texttt{std::tanhf};
допуск~--- относительная погрешность $10^{-4}$.
Тест верифицирует трансцендентные инструкции PTX (\texttt{ex2},
\texttt{lg2}, \texttt{fma}) и преобразования типов.

\textbf{Тестирование.}
Фиксируются: максимальная относительная ошибка, число инструкций
PTX в теле ядра (\texttt{Module::Dump()}), наличие дивергенции
(ожидается $\eta_{\mathrm{branch}} = 1.0$~--- все пути конвергентны).

\textbf{Вывод.}
Верифицирует корректность IEEE~754 для трансцендентных функций;
демонстрирует линейный рост времени с числом инструкций на поток.

\subsubsection{\texttt{softmax\_custom} --- онлайн-нормализация Softmax}

\textbf{Описание теста.}
Ядро реализует трёхфазный онлайн-алгоритм Softmax: поиск максимума,
суммирование экспонент, нормализация. Размер: 1 блок × 32~потока,
длина вектора = 32~элемента. CPU-эталон~--- двухпроходная реализация;
допуск $10^{-4}$.

Тест чувствителен к корректности \texttt{bar.sync} между фазами и
предикатным инструкциям при редукции.

\textbf{Тестирование.}
Проверяются: код возврата, сумма выходного вектора $\approx 1.0$
(свойство Softmax), максимальная ошибка по элементам. Через
профилировщик: число вызовов \texttt{bar.sync} (ожидается~2)
и $\overline{\eta}_{\mathrm{branch}}$ инструкций редукции.

\textbf{Вывод.}
Подтверждает корректность барьерной синхронизации и операций
с разделяемой памятью в многофазных редукциях.

\subsubsection{\texttt{mmul\_custom} --- тайловое матричное умножение}

\textbf{Описание теста.}
Умножение матриц $A(M \times K) \times B(K \times N) = C(M \times N)$
с тайловым алгоритмом (тайл $16 \times 16$). Размер задачи:
$M = N = K = 64$, 16~блоков по 256~потоков. CPU-эталон~---
наивный тройной цикл; допуск~--- относительная погрешность $10^{-3}$.

\textbf{Тестирование.}
Через профилировщик: число уникальных PC в ядре,
$\overline{C}_{\mathrm{bank}}$ для \texttt{ld.shared}
(при правильном заполнении ожидается~0).

\begin{table}[h]
  \caption{Время исполнения \texttt{mmul\_custom} (медиана по 5 повторениям, мс)}
  \label{tab:mmul_perf}
  \centering
  \begin{tabular}{rrrr}
    \toprule
    $M = N = K$ & Без проф., мс & Со сбором проф., мс & Накл. расходы, \% \\
    \midrule
    32  & 4.36   & 28.98   & 565 \\
    64  & 24.22  & 191.87  & 692 \\
    96  & 69.57  & 606.19  & 771 \\
    128 & 146.97 & 1375.96 & 836 \\
    \bottomrule
  \end{tabular}
\end{table}

Высокие накладные расходы (565--836\%) объясняются инструкционной плотностью ядра:
каждый поток выполняет $K$ итераций скалярного произведения, генерируя $\Theta(K)$
PTX-инструкций. Профилировщик создаёт запись \texttt{ProfilingRecord}
для каждого исполнения, что линейно по числу инструкций. При $K=128$
число записей в буфер пропорционально $M^3/\texttt{WARP\_SIZE}$,
откуда следует рост накладных расходов с увеличением размера задачи.

\textbf{Вывод.}
Наиболее ресурсоёмкий тест~--- верхняя граница оценки
производительности эмулятора. Ожидается $O(N^3)$-рост числа инструкций.

\subsubsection{\texttt{attention\_custom} --- Flash Attention-подобное ядро}

\textbf{Описание теста.}
Ядро реализует:
\[
  O = \mathrm{softmax}\!\left(\frac{Q K^\top}{\sqrt{d}}\right) V
\]
с разбивкой на тайлы по последовательности. Структура управляющего
потока: вложенные циклы по тайлам, несколько барьеров, смешанные
обращения к разделяемой и глобальной памяти.
Верификация против наивной CPU-реализации; допуск $10^{-3}$.

\textbf{Тестирование.}
Анализируется: глубина стека дивергенции (ожидается $k \leq 2$),
число вызовов \texttt{bar.sync} на блок,
$\overline{\eta}_{\mathrm{branch}}$ внутри тела тайлового цикла.

\textbf{Вывод.}
Успешное прохождение подтверждает корректность полного стека:
перехват, анализатор, исполнители, PDOM, барьер, разделяемая память.
Является главным системным тестом.

\subsubsection{\texttt{fft\_custom} --- Radix-2 DIT FFT, $N=32$}

\textbf{Описание теста.}
Быстрое преобразование Фурье (алгоритм Cooley-Tukey), $N=32$,
32~серии. Каждый блок из 32~потоков обрабатывает одну серию:
перестановка с обращением битов при загрузке, 5~стадий butterfly,
барьеры между стадиями.
Верификация против CPU-DFT:
$\mathit{rel\_err}(k) \leq 10^{-3}$ для всех частот и серий.

\textbf{Тестирование.}
Из профилировщика: число инструкций в одном проходе butterfly,
$\overline{\eta}_{\mathrm{branch}}$ на операции с разделяемой памятью.

\textbf{Вывод.}
Подтверждает корректность PDOM при 5~последовательных
\texttt{\_\_syncthreads}; устранение аварийного завершения анализатора
при PTX-файле $\approx$ 35~КБ.

\subsubsection{\texttt{conjugate\_gradient} --- итерационное решение СЛАУ}

\textbf{Описание теста.}
Метод сопряжённых градиентов для $Ax = b$. Единственный тест
с \textit{динамически изменяемым} числом итераций: завершение
при $\|r\| < \varepsilon$.
Верификация: $\|Ax^* - b\| / \|b\| \leq 10^{-4}$.

\textbf{Тестирование.}
Из профилировщика: общий \texttt{exec\_count} суммирующих инструкций
(меняется от запуска к запуску), $\overline{\eta}_{\mathrm{branch}}$
условия выхода.

\textbf{Вывод.}
Единственный тест с нефиксированным числом итераций; подтверждает
корректность для итерационных алгоритмов с динамической дивергенцией.

\subsubsection{Сводная таблица результатов интеграционных тестов}

\begin{table}[h]
  \caption{Корректность интеграционных тестов}
  \label{tab:integration_correctness}
  \centering
  \begin{tabular}{lcrr}
    \toprule
    \textbf{Ядро} & \textbf{Статус} & \textbf{Макс. ошибка} & \textbf{PTX инстр.} \\
    \midrule
    \texttt{vadd}       & PASS & $< 10^{-6}$ & 21 \\
    \texttt{gelu}       & PASS & $< 10^{-4}$ & 45 \\
    \texttt{softmax}    & PASS & $< 10^{-4}$ & 54 \\
    \texttt{mmul}       & PASS & $< 10^{-3}$ & 77 \\
    \texttt{attention}  & PASS & $< 10^{-3}$ & 129 \\
    \texttt{fft}        & PASS & $< 10^{-3}$ & 219 \\
    \texttt{conj\_grad} & PASS & $< 10^{-4}$ & 187 \\
    \bottomrule
  \end{tabular}
\end{table}

\begin{table}[h]
  \caption{Производительность на входных данных $\approx 64$~КБ. Сравнение с реальными GPU (V100-PCIE)}
  \label{tab:integration_results}
  \centering
  \footnotesize
  \begin{tabular}{lrrrrrrr}
    \toprule
    \textbf{Ядро} &
    \textbf{Без проф.} & \textbf{Со сбором} & \textbf{НР} &
    \textbf{AI} &
    \textbf{V100} &
    \textbf{Замедл.} \\
    & \textbf{(мс)} & \textbf{(мс)} & \textbf{(\%)} & \textbf{(Ф/Б)} &
    \textbf{(мкс)} & \textbf{(×)} \\
    \midrule
    \texttt{vadd}~(N=8192)       &   20.5 &   34.8 &  70 & 0.17 &   7.2 &  $\approx 2\,800$ \\
    \texttt{gelu}~(N=16384)      &   46.8 &  121.8 & 160 & 1.25 &   8.0 &  $\approx 5\,900$ \\
    \texttt{softmax}~(512×32)    &   60.9 &  268.4 & 341 & 0.38 &   7.8 &  $\approx 7\,800$ \\
    \texttt{mmul}~(128×128)      &  111.7 &  970.5 & 769 & 21.3 &  17.8 &  $\approx 6\,300$ \\
    \texttt{attention}~(S=D=128) &  300.5 & 2074.1 & 590 & 25.8 & 131.9 &  $\approx 2\,300$ \\
    \texttt{fft}~(N=32,B=256)    &   52.6 &  218.0 & 315 & 1.56 &   8.9 &  $\approx 5\,900$ \\
    \texttt{conj\_grad}~(N=4096) &  890.6 & 1879.7 & 111 & 0.47 & 476.3 &  $\approx 1\,900$ \\
    \bottomrule
  \end{tabular}
  \begin{flushleft}
  {\footnotesize НР~--- накладные расходы профилировщика (\%);
  AI~--- арифметическая интенсивность (Флоп/Байт).}
  \end{flushleft}
\end{table}

По результатам тестирования эмулятор корректно воспроизводит вычислительные
результаты всех семи алгоритмически различных ядер при заданных допусках.

Совокупность показателей подтверждает, что реализованный эмулятор
обеспечивает не только синтаксическую корректность разбора PTX,
но и семантически точное исполнение: результаты всех ядер численно
совпадают с CPU-эталонами при жёстких порогах погрешности $10^{-3}$--$10^{-6}$.

\textbf{Сравнение с V100-PCIE.}
Для ядер с нетривиальным временем на V100 (\texttt{attention}~--- 132~мкс,
\texttt{conj\_grad}~--- 476~мкс) замедление сходится к 1\,900--2\,300$\times$~---
это реальная стоимость интерпретируемого однопоточного исполнения PTX
без JIT-компиляции и без параллелизма между варпами.
Для лёгких ядер (\texttt{vadd}, \texttt{gelu}, \texttt{softmax}, \texttt{fft})
замедление составляет 5\,900--7\,800$\times$: V100 исполняет такие ядра
за 7--9~мкс вне зависимости от объёма вычислений, поскольку доминирует
задержка запуска ядра ($\approx$7~мкс на V100-PCIE), тогда как эмулятор
тратит десятки миллисекунд на интерпретацию всех инструкций.
Итоговое отношение является артефактом малого абсолютного времени GPU,
а не показателем превышения ожидаемого замедления.

\subsection{Часть II --- Тестирование системы профилирования}

\textbf{Методология.} Каждый тест профилировщика запускается с флагом
\texttt{--collect-profiling}; полученный \texttt{profiling.txt}
подаётся в \texttt{python -m tools.ptx\_report}; HTML-отчёт верифицируется
против аналитически ожидаемых значений.

Верификация профилировщика производится отдельно от функциональных
интеграционных тестов: ядро может выдавать численно корректный результат
при одновременно ошибочном подсчёте метрик. Раздельные тесты
исключают эту возможность, проверяя каждую метрику независимо
на простых ядрах с заранее известными аналитическими значениями.

\subsubsection{Тест метрики \texttt{branch\_efficiency}}

\textbf{Описание.}
Ядро \texttt{branch\_divergence.cu}: нечётные потоки
($\mathit{tid} \bmod 2 \neq 0$) исполняют одну ветку, чётные~--- другую.
Аналитически: $\eta_{\mathrm{branch}} = 16/32 = 0.5$ для инструкций
внутри каждой ветки.

\textbf{Замер.}
Из \texttt{profiling.txt} фильтруются записи инструкций с предикатным
префиксом. Строится таблица:
\texttt{PC | маска | popcount | $\eta_{\mathrm{emu}}$ | $\eta_{\mathrm{analytic}}$ | $\Delta$}.

\textbf{Вывод.}
Отклонение $\Delta = 0$ (значение детерминировано). Подтверждается
корректность формулы $\eta_{\mathrm{branch}} = \popcount(M)/32$ и
правильность записи маски в \texttt{ProfilingRecord}.

\subsubsection{Тест метрики \texttt{bank\_conflicts}}

\textbf{Описание.}
Ядро \texttt{bank\_conflicts.cu} с тремя шаблонами доступа
к разделяемой памяти (конфиг: $B = 32$, $w_b = 4$ байта):
\begin{enumerate}
  \item Последовательный шаг 4~байта: $C_{\mathrm{bank}} = 0$.
  \item Шаг $w_b \cdot B = 128$~байт: все потоки в банк~0,
        $C_{\mathrm{bank}} = 31$.
  \item Шаг 8~байт: конфликт 2~степени, $C_{\mathrm{bank}} = 1$.
\end{enumerate}

\begin{table}[h]
  \caption{Результаты теста \texttt{bank\_conflicts}}
  \label{tab:bank_conflicts_test}
  \centering
  \begin{tabular}{lrrr}
    \toprule
    \textbf{Шаблон} & $C_{\mathrm{bank,expected}}$ & $C_{\mathrm{bank,emu}}$ & Статус \\
    \midrule
    Шаг 4 B    & 0  & 0  & PASS \\
    Шаг 128 B  & 31 & 31 & PASS \\
    Шаг 8 B    & 1  & 1  & PASS \\
    \bottomrule
  \end{tabular}
\end{table}

\subsubsection{Тест метрики \texttt{global\_coalescing}}

\textbf{Описание.}
Пять шаблонов доступа для \texttt{ld.global.f32} ($s=4$~байта):

\begin{table}[h]
  \caption{Результаты теста \texttt{global\_coalescing}}
  \label{tab:coalescing_test}
  \centering
  \begin{tabular}{lrrrrrr}
    \toprule
    \textbf{Шаблон} & $s$, Б & $T_{\mathrm{ideal}}$ & $\kappa_{\mathrm{exp}}$ &
    $T_{\mathrm{actual,emu}}$ & $\kappa_{\mathrm{emu}}$ & $\Delta$ \\
    \midrule
    Последовательный, f32 & 4 & 1 & 1.000 & 1  & 1.000 & 0.000 \\
    Шаг 128~Б, f32        & 4 & 1 & 0.031 & 32 & 0.031 & 0.000 \\
    Шаг 16~Б, f32         & 4 & 1 & 0.250 & 4  & 0.250 & 0.000 \\
    Широковещательный     & 4 & 1 & 1.000 & 1  & 1.000 & 0.000 \\
    Последовательный, f64 & 8 & 2 & 1.000 & 2  & 1.000 & 0.000 \\
    \bottomrule
  \end{tabular}
\end{table}

\textbf{Вывод.}
Корректность случая \texttt{f64} ($T_{\mathrm{ideal}} = 2$) подтверждает
учёт размера элемента в формуле $T_{\mathrm{ideal}} = \lceil \mathit{active} \cdot s / 128 \rceil$.

\subsubsection{Тест статистики регистров (\texttt{write\_ops})}

\textbf{Описание.}
Ядро \texttt{vadd\_custom}: 3~инструкции с результатом, 1~блок, 1~варп,
полная маска. Аналитически: \texttt{write\_ops[\%f] = 1 × 32 × 1 × 1 = 32}.

\textbf{Вывод.}
Подтверждает корректность детектирования \texttt{dst\_reg\_prefix}
и накопления \texttt{write\_ops}.

\subsubsection{Тест аннотации исходного кода через \texttt{.loc}}

\newpage

\textbf{Описание.}
Ядро \texttt{fft\_custom.cu} компилируется с \texttt{nvcc -lineinfo}.
В HTML-отчёте каждая PTX-инструкция аннотирована строкой
\texttt{fft\_custom.cu:L}. Выборочно проверяются 5~инструкций
с известными \texttt{.loc}-значениями.

\textbf{Вывод.}
Подтверждает корректность разбора \texttt{.file}/\texttt{.loc}
в \texttt{ptx\_parser.py}.

\subsubsection{Тест агрегации: ручной расчёт vs. отчёт}

\textbf{Описание.}
Ядро \texttt{vadd\_custom} с полностью известной структурой.
Аналитически: \texttt{exec\_count = 3~×~32 = 96},
$\overline{\eta}_{\mathrm{branch}} = 1.0$, $C_{\mathrm{bank}} = 0$,
\texttt{global\_mem\_transactions~= 2}.

\begin{table}[h]
  \caption{Результаты теста агрегации}
  \label{tab:aggregation_test}
  \centering
  \begin{tabular}{lrrr}
    \toprule
    \textbf{Метрика} & \textbf{Ожидаемое} & \textbf{HTML-отчёт} & $\Delta$ \\
    \midrule
    \texttt{exec\_count}          & 96    & 96    & 0 \\
    $\overline{\eta}_{\mathrm{branch}}$  & 1.000 & 1.000 & 0.000 \\
    \texttt{bank\_conflicts\_total} & 0     & 0     & 0 \\
    \texttt{global\_mem\_txns}     & 2     & 2     & 0 \\
    \bottomrule
  \end{tabular}
\end{table}

\textbf{Вывод.}
Нулевые отклонения подтверждают корректность всей цепочки:
\texttt{profiling.txt} $\to$ \texttt{parser.py} $\to$
\texttt{aggregator.py} $\to$ \texttt{renderer.py}.

\newpage
